% zuka-pro.tex — última capa de la plantilla (2026-11-24, spec 2026-11-24-libros-profesional).
% Estilo editorial sobrio: tinta azul noche de la marca, un solo acento (lima) y líneas finas en lugar de cajas de colores.
% En Razonamiento lógico cada lección abre con la ruta de la explicación de la app (explora-rl.tex, generado) y cierra con
% «Lo que construiste»; piezas \capitulo, \turno, \armar.
% GENERADO por docs/libros/arte/desde_app.js: no editar a mano.
% \zukamark[color]{alto}: el cerebro; \zukahor: cerebro y palabra en fila; \zukaver: la palabra debajo.
\newcommand{\zuka@dibmark}{\pgfpathmoveto{\pgfpoint{49.522mm}{-71.730mm}}%
\pgfpathcurveto{\pgfpoint{48.604mm}{-71.523mm}}{\pgfpoint{48.046mm}{-71.029mm}}{\pgfpoint{47.592mm}{-70.008mm}}%
\pgfpathcurveto{\pgfpoint{47.209mm}{-69.171mm}}{\pgfpoint{46.172mm}{-67.026mm}}{\pgfpoint{45.646mm}{-66.013mm}}%
\pgfpathcurveto{\pgfpoint{45.064mm}{-64.872mm}}{\pgfpoint{44.793mm}{-64.466mm}}{\pgfpoint{44.274mm}{-63.955mm}}%
\pgfpathcurveto{\pgfpoint{43.700mm}{-63.373mm}}{\pgfpoint{43.573mm}{-63.349mm}}{\pgfpoint{42.903mm}{-63.652mm}}%
\pgfpathcurveto{\pgfpoint{38.142mm}{-65.837mm}}{\pgfpoint{32.767mm}{-63.509mm}}{\pgfpoint{31.300mm}{-58.636mm}}%
\pgfpathcurveto{\pgfpoint{31.108mm}{-57.990mm}}{\pgfpoint{31.132mm}{-58.014mm}}{\pgfpoint{30.582mm}{-57.935mm}}%
\pgfpathcurveto{\pgfpoint{26.563mm}{-57.337mm}}{\pgfpoint{23.333mm}{-54.665mm}}{\pgfpoint{22.002mm}{-50.821mm}}%
\pgfpathcurveto{\pgfpoint{20.662mm}{-46.978mm}}{\pgfpoint{21.595mm}{-43.006mm}}{\pgfpoint{24.545mm}{-40.048mm}}%
\pgfpathcurveto{\pgfpoint{25.056mm}{-39.530mm}}{\pgfpoint{25.048mm}{-39.569mm}}{\pgfpoint{24.817mm}{-39.099mm}}%
\pgfpathcurveto{\pgfpoint{23.030mm}{-35.526mm}}{\pgfpoint{23.836mm}{-31.005mm}}{\pgfpoint{26.802mm}{-27.943mm}}%
\pgfpathcurveto{\pgfpoint{28.405mm}{-26.284mm}}{\pgfpoint{30.359mm}{-25.295mm}}{\pgfpoint{32.592mm}{-25.016mm}}%
\pgfpathcurveto{\pgfpoint{33.038mm}{-24.960mm}}{\pgfpoint{33.054mm}{-24.936mm}}{\pgfpoint{33.094mm}{-24.522mm}}%
\pgfpathcurveto{\pgfpoint{33.293mm}{-22.448mm}}{\pgfpoint{34.091mm}{-20.670mm}}{\pgfpoint{35.558mm}{-19.059mm}}%
\pgfpathcurveto{\pgfpoint{39.282mm}{-14.952mm}}{\pgfpoint{45.813mm}{-14.649mm}}{\pgfpoint{49.641mm}{-18.397mm}}%
\pgfpathcurveto{\pgfpoint{50.199mm}{-18.939mm}}{\pgfpoint{50.152mm}{-18.939mm}}{\pgfpoint{50.566mm}{-18.525mm}}%
\pgfpathcurveto{\pgfpoint{53.868mm}{-15.215mm}}{\pgfpoint{59.155mm}{-14.841mm}}{\pgfpoint{63.078mm}{-17.640mm}}%
\pgfpathcurveto{\pgfpoint{65.359mm}{-19.266mm}}{\pgfpoint{66.707mm}{-21.619mm}}{\pgfpoint{66.962mm}{-24.426mm}}%
\pgfpathcurveto{\pgfpoint{67.010mm}{-24.944mm}}{\pgfpoint{67.026mm}{-24.960mm}}{\pgfpoint{67.424mm}{-25.016mm}}%
\pgfpathcurveto{\pgfpoint{70.949mm}{-25.447mm}}{\pgfpoint{74.219mm}{-28.086mm}}{\pgfpoint{75.463mm}{-31.499mm}}%
\pgfpathcurveto{\pgfpoint{76.236mm}{-33.612mm}}{\pgfpoint{76.244mm}{-36.244mm}}{\pgfpoint{75.494mm}{-38.238mm}}%
\pgfpathcurveto{\pgfpoint{75.375mm}{-38.541mm}}{\pgfpoint{75.136mm}{-39.099mm}}{\pgfpoint{74.984mm}{-39.402mm}}%
\pgfpathcurveto{\pgfpoint{74.809mm}{-39.745mm}}{\pgfpoint{74.809mm}{-39.745mm}}{\pgfpoint{75.343mm}{-40.295mm}}%
\pgfpathcurveto{\pgfpoint{76.810mm}{-41.826mm}}{\pgfpoint{77.719mm}{-43.676mm}}{\pgfpoint{78.030mm}{-45.774mm}}%
\pgfpathcurveto{\pgfpoint{78.804mm}{-51.077mm}}{\pgfpoint{75.431mm}{-56.164mm}}{\pgfpoint{70.263mm}{-57.480mm}}%
\pgfpathcurveto{\pgfpoint{69.825mm}{-57.592mm}}{\pgfpoint{69.609mm}{-57.632mm}}{\pgfpoint{68.971mm}{-57.719mm}}%
\pgfpathcurveto{\pgfpoint{68.525mm}{-57.783mm}}{\pgfpoint{68.549mm}{-57.759mm}}{\pgfpoint{68.365mm}{-58.453mm}}%
\pgfpathcurveto{\pgfpoint{67.217mm}{-62.903mm}}{\pgfpoint{62.687mm}{-65.415mm}}{\pgfpoint{58.230mm}{-64.075mm}}%
\pgfpathcurveto{\pgfpoint{57.368mm}{-63.812mm}}{\pgfpoint{57.049mm}{-64.035mm}}{\pgfpoint{55.869mm}{-65.710mm}}%
\pgfpathcurveto{\pgfpoint{55.048mm}{-66.874mm}}{\pgfpoint{53.150mm}{-69.442mm}}{\pgfpoint{52.241mm}{-70.630mm}}%
\pgfpathcurveto{\pgfpoint{51.555mm}{-71.507mm}}{\pgfpoint{50.470mm}{-71.954mm}}{\pgfpoint{49.522mm}{-71.730mm}}%
\pgfpathclose%
\pgfpathmoveto{\pgfpoint{51.826mm}{-64.522mm}}%
\pgfpathcurveto{\pgfpoint{51.922mm}{-64.442mm}}{\pgfpoint{52.520mm}{-63.596mm}}{\pgfpoint{53.437mm}{-62.241mm}}%
\pgfpathcurveto{\pgfpoint{54.992mm}{-59.952mm}}{\pgfpoint{56.914mm}{-59.123mm}}{\pgfpoint{59.043mm}{-59.817mm}}%
\pgfpathcurveto{\pgfpoint{59.195mm}{-59.872mm}}{\pgfpoint{59.458mm}{-59.960mm}}{\pgfpoint{59.633mm}{-60.024mm}}%
\pgfpathcurveto{\pgfpoint{61.467mm}{-60.662mm}}{\pgfpoint{63.565mm}{-59.625mm}}{\pgfpoint{64.314mm}{-57.703mm}}%
\pgfpathcurveto{\pgfpoint{65.287mm}{-55.231mm}}{\pgfpoint{63.549mm}{-52.648mm}}{\pgfpoint{60.805mm}{-52.496mm}}%
\pgfpathcurveto{\pgfpoint{59.545mm}{-52.432mm}}{\pgfpoint{58.876mm}{-51.659mm}}{\pgfpoint{59.075mm}{-50.518mm}}%
\pgfpathcurveto{\pgfpoint{59.242mm}{-49.577mm}}{\pgfpoint{59.944mm}{-49.155mm}}{\pgfpoint{61.188mm}{-49.250mm}}%
\pgfpathcurveto{\pgfpoint{63.915mm}{-49.442mm}}{\pgfpoint{66.196mm}{-50.869mm}}{\pgfpoint{67.456mm}{-53.190mm}}%
\pgfpathcurveto{\pgfpoint{67.727mm}{-53.684mm}}{\pgfpoint{67.727mm}{-53.684mm}}{\pgfpoint{68.126mm}{-53.676mm}}%
\pgfpathcurveto{\pgfpoint{70.375mm}{-53.652mm}}{\pgfpoint{72.903mm}{-51.627mm}}{\pgfpoint{73.692mm}{-49.203mm}}%
\pgfpathcurveto{\pgfpoint{75.271mm}{-44.402mm}}{\pgfpoint{70.566mm}{-39.817mm}}{\pgfpoint{65.877mm}{-41.555mm}}%
\pgfpathcurveto{\pgfpoint{64.290mm}{-42.145mm}}{\pgfpoint{62.951mm}{-40.805mm}}{\pgfpoint{63.652mm}{-39.346mm}}%
\pgfpathcurveto{\pgfpoint{63.907mm}{-38.812mm}}{\pgfpoint{64.179mm}{-38.620mm}}{\pgfpoint{65.104mm}{-38.301mm}}%
\pgfpathcurveto{\pgfpoint{66.834mm}{-37.719mm}}{\pgfpoint{68.341mm}{-37.648mm}}{\pgfpoint{70.136mm}{-38.062mm}}%
\pgfpathcurveto{\pgfpoint{70.893mm}{-38.246mm}}{\pgfpoint{70.917mm}{-38.230mm}}{\pgfpoint{71.340mm}{-37.400mm}}%
\pgfpathcurveto{\pgfpoint{72.871mm}{-34.418mm}}{\pgfpoint{71.324mm}{-30.590mm}}{\pgfpoint{68.046mm}{-29.282mm}}%
\pgfpathcurveto{\pgfpoint{65.941mm}{-28.437mm}}{\pgfpoint{63.501mm}{-29.075mm}}{\pgfpoint{62.217mm}{-30.821mm}}%
\pgfpathcurveto{\pgfpoint{61.077mm}{-32.368mm}}{\pgfpoint{58.756mm}{-31.396mm}}{\pgfpoint{59.187mm}{-29.561mm}}%
\pgfpathcurveto{\pgfpoint{59.250mm}{-29.282mm}}{\pgfpoint{59.386mm}{-29.043mm}}{\pgfpoint{59.665mm}{-28.708mm}}%
\pgfpathcurveto{\pgfpoint{60.646mm}{-27.536mm}}{\pgfpoint{61.252mm}{-27.018mm}}{\pgfpoint{62.352mm}{-26.404mm}}%
\pgfpathcurveto{\pgfpoint{62.823mm}{-26.148mm}}{\pgfpoint{62.839mm}{-26.108mm}}{\pgfpoint{62.847mm}{-25.359mm}}%
\pgfpathcurveto{\pgfpoint{62.847mm}{-20.311mm}}{\pgfpoint{56.332mm}{-17.982mm}}{\pgfpoint{52.919mm}{-21.810mm}}%
\pgfpathcurveto{\pgfpoint{51.962mm}{-22.879mm}}{\pgfpoint{51.627mm}{-23.876mm}}{\pgfpoint{51.603mm}{-25.734mm}}%
\pgfpathcurveto{\pgfpoint{51.587mm}{-26.746mm}}{\pgfpoint{51.579mm}{-26.802mm}}{\pgfpoint{51.411mm}{-27.129mm}}%
\pgfpathcurveto{\pgfpoint{50.797mm}{-28.317mm}}{\pgfpoint{49.003mm}{-28.166mm}}{\pgfpoint{48.541mm}{-26.890mm}}%
\pgfpathcurveto{\pgfpoint{48.477mm}{-26.707mm}}{\pgfpoint{48.477mm}{-26.691mm}}{\pgfpoint{48.453mm}{-25.622mm}}%
\pgfpathcurveto{\pgfpoint{48.437mm}{-24.482mm}}{\pgfpoint{48.421mm}{-24.338mm}}{\pgfpoint{48.285mm}{-23.820mm}}%
\pgfpathcurveto{\pgfpoint{47.305mm}{-20.215mm}}{\pgfpoint{42.823mm}{-18.740mm}}{\pgfpoint{39.530mm}{-20.941mm}}%
\pgfpathcurveto{\pgfpoint{37.943mm}{-22.002mm}}{\pgfpoint{36.858mm}{-23.955mm}}{\pgfpoint{36.962mm}{-25.558mm}}%
\pgfpathcurveto{\pgfpoint{36.986mm}{-26.045mm}}{\pgfpoint{37.010mm}{-26.069mm}}{\pgfpoint{37.504mm}{-26.316mm}}%
\pgfpathcurveto{\pgfpoint{38.987mm}{-27.049mm}}{\pgfpoint{40.686mm}{-28.676mm}}{\pgfpoint{40.893mm}{-29.561mm}}%
\pgfpathcurveto{\pgfpoint{41.332mm}{-31.435mm}}{\pgfpoint{38.987mm}{-32.376mm}}{\pgfpoint{37.831mm}{-30.797mm}}%
\pgfpathcurveto{\pgfpoint{36.874mm}{-29.482mm}}{\pgfpoint{34.944mm}{-28.684mm}}{\pgfpoint{33.262mm}{-28.907mm}}%
\pgfpathcurveto{\pgfpoint{29.250mm}{-29.434mm}}{\pgfpoint{26.730mm}{-33.876mm}}{\pgfpoint{28.469mm}{-37.360mm}}%
\pgfpathcurveto{\pgfpoint{28.876mm}{-38.182mm}}{\pgfpoint{28.907mm}{-38.198mm}}{\pgfpoint{29.537mm}{-38.046mm}}%
\pgfpathcurveto{\pgfpoint{31.443mm}{-37.584mm}}{\pgfpoint{33.086mm}{-37.648mm}}{\pgfpoint{34.888mm}{-38.254mm}}%
\pgfpathcurveto{\pgfpoint{36.061mm}{-38.652mm}}{\pgfpoint{36.515mm}{-39.187mm}}{\pgfpoint{36.459mm}{-40.112mm}}%
\pgfpathcurveto{\pgfpoint{36.388mm}{-41.332mm}}{\pgfpoint{35.303mm}{-41.930mm}}{\pgfpoint{34.027mm}{-41.451mm}}%
\pgfpathcurveto{\pgfpoint{30.008mm}{-39.960mm}}{\pgfpoint{25.463mm}{-43.094mm}}{\pgfpoint{25.463mm}{-47.368mm}}%
\pgfpathcurveto{\pgfpoint{25.463mm}{-50.582mm}}{\pgfpoint{28.118mm}{-53.525mm}}{\pgfpoint{31.300mm}{-53.852mm}}%
\pgfpathcurveto{\pgfpoint{31.850mm}{-53.907mm}}{\pgfpoint{31.946mm}{-53.876mm}}{\pgfpoint{32.121mm}{-53.501mm}}%
\pgfpathcurveto{\pgfpoint{33.373mm}{-50.861mm}}{\pgfpoint{36.093mm}{-49.179mm}}{\pgfpoint{39.091mm}{-49.179mm}}%
\pgfpathcurveto{\pgfpoint{40.112mm}{-49.179mm}}{\pgfpoint{40.774mm}{-49.649mm}}{\pgfpoint{40.925mm}{-50.494mm}}%
\pgfpathcurveto{\pgfpoint{41.124mm}{-51.547mm}}{\pgfpoint{40.407mm}{-52.416mm}}{\pgfpoint{39.338mm}{-52.416mm}}%
\pgfpathcurveto{\pgfpoint{37.711mm}{-52.416mm}}{\pgfpoint{36.427mm}{-53.062mm}}{\pgfpoint{35.630mm}{-54.282mm}}%
\pgfpathcurveto{\pgfpoint{33.892mm}{-56.922mm}}{\pgfpoint{35.925mm}{-60.263mm}}{\pgfpoint{39.258mm}{-60.263mm}}%
\pgfpathcurveto{\pgfpoint{40.064mm}{-60.263mm}}{\pgfpoint{40.534mm}{-60.128mm}}{\pgfpoint{41.531mm}{-59.633mm}}%
\pgfpathcurveto{\pgfpoint{43.931mm}{-58.445mm}}{\pgfpoint{45.781mm}{-59.115mm}}{\pgfpoint{47.895mm}{-61.946mm}}%
\pgfpathcurveto{\pgfpoint{48.285mm}{-62.464mm}}{\pgfpoint{48.341mm}{-62.512mm}}{\pgfpoint{48.413mm}{-62.376mm}}%
\pgfpathcurveto{\pgfpoint{48.437mm}{-62.337mm}}{\pgfpoint{48.445mm}{-61.531mm}}{\pgfpoint{48.461mm}{-58.437mm}}%
\pgfpathcurveto{\pgfpoint{48.469mm}{-54.561mm}}{\pgfpoint{48.469mm}{-54.561mm}}{\pgfpoint{48.525mm}{-54.410mm}}%
\pgfpathcurveto{\pgfpoint{49.075mm}{-52.982mm}}{\pgfpoint{51.005mm}{-52.982mm}}{\pgfpoint{51.523mm}{-54.410mm}}%
\pgfpathcurveto{\pgfpoint{51.579mm}{-54.585mm}}{\pgfpoint{51.579mm}{-54.585mm}}{\pgfpoint{51.595mm}{-59.530mm}}%
\pgfpathcurveto{\pgfpoint{51.611mm}{-64.793mm}}{\pgfpoint{51.603mm}{-64.569mm}}{\pgfpoint{51.730mm}{-64.569mm}}%
\pgfpathcurveto{\pgfpoint{51.754mm}{-64.569mm}}{\pgfpoint{51.802mm}{-64.545mm}}{\pgfpoint{51.826mm}{-64.522mm}}%
\pgfpathclose\pgfusepath{fill}%
\pgfpathmoveto{\pgfpoint{67.903mm}{-18.445mm}}%
\pgfpathcurveto{\pgfpoint{67.033mm}{-18.317mm}}{\pgfpoint{66.427mm}{-17.552mm}}{\pgfpoint{66.491mm}{-16.667mm}}%
\pgfpathcurveto{\pgfpoint{66.507mm}{-16.356mm}}{\pgfpoint{66.555mm}{-16.228mm}}{\pgfpoint{67.018mm}{-15.167mm}}%
\pgfpathcurveto{\pgfpoint{67.233mm}{-14.657mm}}{\pgfpoint{67.544mm}{-13.931mm}}{\pgfpoint{67.703mm}{-13.557mm}}%
\pgfpathcurveto{\pgfpoint{67.871mm}{-13.182mm}}{\pgfpoint{68.078mm}{-12.687mm}}{\pgfpoint{68.182mm}{-12.456mm}}%
\pgfpathcurveto{\pgfpoint{68.285mm}{-12.217mm}}{\pgfpoint{68.493mm}{-11.730mm}}{\pgfpoint{68.652mm}{-11.364mm}}%
\pgfpathcurveto{\pgfpoint{69.123mm}{-10.247mm}}{\pgfpoint{69.131mm}{-10.231mm}}{\pgfpoint{69.386mm}{-9.984mm}}%
\pgfpathcurveto{\pgfpoint{70.486mm}{-8.892mm}}{\pgfpoint{72.305mm}{-9.633mm}}{\pgfpoint{72.305mm}{-11.172mm}}%
\pgfpathcurveto{\pgfpoint{72.305mm}{-11.499mm}}{\pgfpoint{72.257mm}{-11.667mm}}{\pgfpoint{72.049mm}{-12.153mm}}%
\pgfpathcurveto{\pgfpoint{71.954mm}{-12.376mm}}{\pgfpoint{71.802mm}{-12.735mm}}{\pgfpoint{71.715mm}{-12.943mm}}%
\pgfpathcurveto{\pgfpoint{71.627mm}{-13.158mm}}{\pgfpoint{71.475mm}{-13.517mm}}{\pgfpoint{71.372mm}{-13.740mm}}%
\pgfpathcurveto{\pgfpoint{71.268mm}{-13.971mm}}{\pgfpoint{71.085mm}{-14.394mm}}{\pgfpoint{70.957mm}{-14.689mm}}%
\pgfpathcurveto{\pgfpoint{70.829mm}{-14.976mm}}{\pgfpoint{70.574mm}{-15.558mm}}{\pgfpoint{70.399mm}{-15.973mm}}%
\pgfpathcurveto{\pgfpoint{70.215mm}{-16.396mm}}{\pgfpoint{69.992mm}{-16.914mm}}{\pgfpoint{69.896mm}{-17.137mm}}%
\pgfpathcurveto{\pgfpoint{69.490mm}{-18.118mm}}{\pgfpoint{68.772mm}{-18.589mm}}{\pgfpoint{67.903mm}{-18.445mm}}%
\pgfpathclose%
\pgfpathmoveto{\pgfpoint{72.129mm}{-23.126mm}}%
\pgfpathcurveto{\pgfpoint{70.837mm}{-22.759mm}}{\pgfpoint{70.431mm}{-21.116mm}}{\pgfpoint{71.404mm}{-20.183mm}}%
\pgfpathcurveto{\pgfpoint{71.531mm}{-20.064mm}}{\pgfpoint{73.174mm}{-18.884mm}}{\pgfpoint{74.346mm}{-18.070mm}}%
\pgfpathcurveto{\pgfpoint{74.514mm}{-17.951mm}}{\pgfpoint{74.745mm}{-17.791mm}}{\pgfpoint{74.856mm}{-17.703mm}}%
\pgfpathcurveto{\pgfpoint{74.960mm}{-17.624mm}}{\pgfpoint{75.247mm}{-17.424mm}}{\pgfpoint{75.494mm}{-17.257mm}}%
\pgfpathcurveto{\pgfpoint{75.734mm}{-17.089mm}}{\pgfpoint{76.085mm}{-16.842mm}}{\pgfpoint{76.276mm}{-16.707mm}}%
\pgfpathcurveto{\pgfpoint{76.810mm}{-16.332mm}}{\pgfpoint{76.986mm}{-16.228mm}}{\pgfpoint{77.129mm}{-16.164mm}}%
\pgfpathcurveto{\pgfpoint{78.517mm}{-15.590mm}}{\pgfpoint{79.944mm}{-16.986mm}}{\pgfpoint{79.442mm}{-18.429mm}}%
\pgfpathcurveto{\pgfpoint{79.290mm}{-18.868mm}}{\pgfpoint{79.075mm}{-19.091mm}}{\pgfpoint{78.246mm}{-19.657mm}}%
\pgfpathcurveto{\pgfpoint{77.911mm}{-19.896mm}}{\pgfpoint{77.329mm}{-20.295mm}}{\pgfpoint{76.962mm}{-20.550mm}}%
\pgfpathcurveto{\pgfpoint{76.595mm}{-20.805mm}}{\pgfpoint{76.085mm}{-21.156mm}}{\pgfpoint{75.829mm}{-21.332mm}}%
\pgfpathcurveto{\pgfpoint{75.582mm}{-21.507mm}}{\pgfpoint{75.144mm}{-21.802mm}}{\pgfpoint{74.864mm}{-21.994mm}}%
\pgfpathcurveto{\pgfpoint{73.971mm}{-22.616mm}}{\pgfpoint{73.461mm}{-22.959mm}}{\pgfpoint{73.341mm}{-23.014mm}}%
\pgfpathcurveto{\pgfpoint{73.030mm}{-23.174mm}}{\pgfpoint{72.464mm}{-23.230mm}}{\pgfpoint{72.129mm}{-23.126mm}}%
\pgfpathclose\pgfusepath{fill}%
}
\newcommand{\zuka@dibword}{\pgfpathmoveto{\pgfpoint{19.689mm}{-87.998mm}}%
\pgfpathcurveto{\pgfpoint{18.828mm}{-87.791mm}}{\pgfpoint{18.469mm}{-86.754mm}}{\pgfpoint{18.995mm}{-85.965mm}}%
\pgfpathcurveto{\pgfpoint{19.059mm}{-85.869mm}}{\pgfpoint{19.338mm}{-85.542mm}}{\pgfpoint{19.625mm}{-85.247mm}}%
\pgfpathcurveto{\pgfpoint{19.904mm}{-84.952mm}}{\pgfpoint{20.223mm}{-84.609mm}}{\pgfpoint{20.335mm}{-84.490mm}}%
\pgfpathcurveto{\pgfpoint{20.439mm}{-84.370mm}}{\pgfpoint{20.734mm}{-84.059mm}}{\pgfpoint{20.981mm}{-83.796mm}}%
\pgfpathcurveto{\pgfpoint{21.236mm}{-83.533mm}}{\pgfpoint{21.587mm}{-83.158mm}}{\pgfpoint{21.770mm}{-82.959mm}}%
\pgfpathcurveto{\pgfpoint{21.954mm}{-82.767mm}}{\pgfpoint{22.352mm}{-82.337mm}}{\pgfpoint{22.663mm}{-82.002mm}}%
\pgfpathcurveto{\pgfpoint{23.517mm}{-81.093mm}}{\pgfpoint{24.147mm}{-80.407mm}}{\pgfpoint{24.450mm}{-80.064mm}}%
\pgfpathcurveto{\pgfpoint{24.601mm}{-79.896mm}}{\pgfpoint{24.896mm}{-79.577mm}}{\pgfpoint{25.104mm}{-79.362mm}}%
\pgfpathcurveto{\pgfpoint{25.534mm}{-78.907mm}}{\pgfpoint{25.598mm}{-78.828mm}}{\pgfpoint{25.598mm}{-78.740mm}}%
\pgfpathcurveto{\pgfpoint{25.598mm}{-78.596mm}}{\pgfpoint{25.678mm}{-78.604mm}}{\pgfpoint{22.759mm}{-78.604mm}}%
\pgfpathcurveto{\pgfpoint{20.534mm}{-78.604mm}}{\pgfpoint{20.056mm}{-78.596mm}}{\pgfpoint{19.904mm}{-78.565mm}}%
\pgfpathcurveto{\pgfpoint{18.644mm}{-78.293mm}}{\pgfpoint{18.636mm}{-76.451mm}}{\pgfpoint{19.888mm}{-76.116mm}}%
\pgfpathcurveto{\pgfpoint{20.112mm}{-76.053mm}}{\pgfpoint{28.812mm}{-76.061mm}}{\pgfpoint{29.043mm}{-76.116mm}}%
\pgfpathcurveto{\pgfpoint{29.657mm}{-76.276mm}}{\pgfpoint{30.016mm}{-76.802mm}}{\pgfpoint{29.976mm}{-77.480mm}}%
\pgfpathcurveto{\pgfpoint{29.952mm}{-77.887mm}}{\pgfpoint{29.841mm}{-78.094mm}}{\pgfpoint{29.386mm}{-78.581mm}}%
\pgfpathcurveto{\pgfpoint{29.043mm}{-78.939mm}}{\pgfpoint{28.421mm}{-79.617mm}}{\pgfpoint{28.070mm}{-80.008mm}}%
\pgfpathcurveto{\pgfpoint{27.935mm}{-80.159mm}}{\pgfpoint{27.632mm}{-80.494mm}}{\pgfpoint{27.408mm}{-80.742mm}}%
\pgfpathcurveto{\pgfpoint{27.177mm}{-80.989mm}}{\pgfpoint{26.786mm}{-81.427mm}}{\pgfpoint{26.531mm}{-81.715mm}}%
\pgfpathcurveto{\pgfpoint{26.276mm}{-81.994mm}}{\pgfpoint{25.853mm}{-82.456mm}}{\pgfpoint{25.590mm}{-82.743mm}}%
\pgfpathcurveto{\pgfpoint{25.327mm}{-83.030mm}}{\pgfpoint{24.984mm}{-83.413mm}}{\pgfpoint{24.825mm}{-83.596mm}}%
\pgfpathcurveto{\pgfpoint{24.657mm}{-83.780mm}}{\pgfpoint{24.330mm}{-84.147mm}}{\pgfpoint{24.083mm}{-84.410mm}}%
\pgfpathcurveto{\pgfpoint{23.301mm}{-85.255mm}}{\pgfpoint{23.293mm}{-85.263mm}}{\pgfpoint{23.325mm}{-85.367mm}}%
\pgfpathcurveto{\pgfpoint{23.349mm}{-85.470mm}}{\pgfpoint{23.373mm}{-85.470mm}}{\pgfpoint{26.411mm}{-85.486mm}}%
\pgfpathcurveto{\pgfpoint{29.386mm}{-85.502mm}}{\pgfpoint{29.386mm}{-85.502mm}}{\pgfpoint{29.593mm}{-85.606mm}}%
\pgfpathcurveto{\pgfpoint{30.694mm}{-86.140mm}}{\pgfpoint{30.494mm}{-87.791mm}}{\pgfpoint{29.306mm}{-88.006mm}}%
\pgfpathcurveto{\pgfpoint{29.051mm}{-88.054mm}}{\pgfpoint{19.888mm}{-88.046mm}}{\pgfpoint{19.689mm}{-87.998mm}}%
\pgfpathclose%
\pgfpathmoveto{\pgfpoint{41.467mm}{-88.525mm}}%
\pgfpathcurveto{\pgfpoint{41.419mm}{-88.525mm}}{\pgfpoint{41.244mm}{-88.501mm}}{\pgfpoint{41.085mm}{-88.493mm}}%
\pgfpathcurveto{\pgfpoint{38.676mm}{-88.262mm}}{\pgfpoint{36.794mm}{-86.611mm}}{\pgfpoint{36.316mm}{-84.306mm}}%
\pgfpathcurveto{\pgfpoint{36.164mm}{-83.596mm}}{\pgfpoint{36.172mm}{-83.780mm}}{\pgfpoint{36.156mm}{-80.423mm}}%
\pgfpathcurveto{\pgfpoint{36.148mm}{-77.033mm}}{\pgfpoint{36.148mm}{-77.105mm}}{\pgfpoint{36.292mm}{-76.778mm}}%
\pgfpathcurveto{\pgfpoint{36.738mm}{-75.805mm}}{\pgfpoint{38.262mm}{-75.805mm}}{\pgfpoint{38.740mm}{-76.770mm}}%
\pgfpathcurveto{\pgfpoint{38.884mm}{-77.073mm}}{\pgfpoint{38.892mm}{-77.089mm}}{\pgfpoint{38.892mm}{-80.159mm}}%
\pgfpathcurveto{\pgfpoint{38.892mm}{-83.134mm}}{\pgfpoint{38.900mm}{-83.357mm}}{\pgfpoint{39.011mm}{-83.812mm}}%
\pgfpathcurveto{\pgfpoint{39.641mm}{-86.292mm}}{\pgfpoint{43.405mm}{-86.635mm}}{\pgfpoint{44.561mm}{-84.314mm}}%
\pgfpathcurveto{\pgfpoint{44.872mm}{-83.700mm}}{\pgfpoint{44.872mm}{-83.748mm}}{\pgfpoint{44.880mm}{-80.215mm}}%
\pgfpathcurveto{\pgfpoint{44.888mm}{-77.488mm}}{\pgfpoint{44.896mm}{-77.145mm}}{\pgfpoint{44.936mm}{-77.010mm}}%
\pgfpathcurveto{\pgfpoint{45.311mm}{-75.718mm}}{\pgfpoint{47.233mm}{-75.734mm}}{\pgfpoint{47.560mm}{-77.033mm}}%
\pgfpathcurveto{\pgfpoint{47.632mm}{-77.329mm}}{\pgfpoint{47.608mm}{-83.413mm}}{\pgfpoint{47.536mm}{-83.868mm}}%
\pgfpathcurveto{\pgfpoint{47.297mm}{-85.247mm}}{\pgfpoint{46.874mm}{-86.116mm}}{\pgfpoint{46.013mm}{-86.970mm}}%
\pgfpathcurveto{\pgfpoint{44.920mm}{-88.046mm}}{\pgfpoint{43.078mm}{-88.676mm}}{\pgfpoint{41.467mm}{-88.525mm}}%
\pgfpathclose%
\pgfpathmoveto{\pgfpoint{55.439mm}{-88.278mm}}%
\pgfpathcurveto{\pgfpoint{54.856mm}{-88.206mm}}{\pgfpoint{54.426mm}{-87.831mm}}{\pgfpoint{54.290mm}{-87.265mm}}%
\pgfpathcurveto{\pgfpoint{54.234mm}{-87.010mm}}{\pgfpoint{54.234mm}{-77.281mm}}{\pgfpoint{54.290mm}{-77.049mm}}%
\pgfpathcurveto{\pgfpoint{54.553mm}{-76.013mm}}{\pgfpoint{55.829mm}{-75.670mm}}{\pgfpoint{56.595mm}{-76.419mm}}%
\pgfpathcurveto{\pgfpoint{56.962mm}{-76.778mm}}{\pgfpoint{56.962mm}{-76.818mm}}{\pgfpoint{56.962mm}{-79.051mm}}%
\pgfpathcurveto{\pgfpoint{56.962mm}{-81.443mm}}{\pgfpoint{56.914mm}{-81.332mm}}{\pgfpoint{57.640mm}{-80.662mm}}%
\pgfpathcurveto{\pgfpoint{58.222mm}{-80.128mm}}{\pgfpoint{59.681mm}{-78.772mm}}{\pgfpoint{60.447mm}{-78.054mm}}%
\pgfpathcurveto{\pgfpoint{62.400mm}{-76.236mm}}{\pgfpoint{62.337mm}{-76.292mm}}{\pgfpoint{62.568mm}{-76.180mm}}%
\pgfpathcurveto{\pgfpoint{63.748mm}{-75.614mm}}{\pgfpoint{64.912mm}{-76.730mm}}{\pgfpoint{64.338mm}{-77.871mm}}%
\pgfpathcurveto{\pgfpoint{64.234mm}{-78.086mm}}{\pgfpoint{63.876mm}{-78.429mm}}{\pgfpoint{62.281mm}{-79.848mm}}%
\pgfpathcurveto{\pgfpoint{60.686mm}{-81.276mm}}{\pgfpoint{60.391mm}{-81.555mm}}{\pgfpoint{60.391mm}{-81.691mm}}%
\pgfpathcurveto{\pgfpoint{60.391mm}{-81.810mm}}{\pgfpoint{60.470mm}{-81.898mm}}{\pgfpoint{62.783mm}{-84.346mm}}%
\pgfpathcurveto{\pgfpoint{63.054mm}{-84.625mm}}{\pgfpoint{63.405mm}{-85.000mm}}{\pgfpoint{63.557mm}{-85.167mm}}%
\pgfpathcurveto{\pgfpoint{63.708mm}{-85.335mm}}{\pgfpoint{63.995mm}{-85.638mm}}{\pgfpoint{64.187mm}{-85.845mm}}%
\pgfpathcurveto{\pgfpoint{64.817mm}{-86.499mm}}{\pgfpoint{64.944mm}{-86.970mm}}{\pgfpoint{64.657mm}{-87.560mm}}%
\pgfpathcurveto{\pgfpoint{64.234mm}{-88.413mm}}{\pgfpoint{63.054mm}{-88.557mm}}{\pgfpoint{62.352mm}{-87.839mm}}%
\pgfpathcurveto{\pgfpoint{62.097mm}{-87.584mm}}{\pgfpoint{61.132mm}{-86.547mm}}{\pgfpoint{60.072mm}{-85.391mm}}%
\pgfpathcurveto{\pgfpoint{59.561mm}{-84.841mm}}{\pgfpoint{58.772mm}{-84.003mm}}{\pgfpoint{58.501mm}{-83.724mm}}%
\pgfpathcurveto{\pgfpoint{58.110mm}{-83.325mm}}{\pgfpoint{58.158mm}{-83.317mm}}{\pgfpoint{57.464mm}{-83.939mm}}%
\pgfpathcurveto{\pgfpoint{56.922mm}{-84.426mm}}{\pgfpoint{56.978mm}{-84.211mm}}{\pgfpoint{56.962mm}{-85.829mm}}%
\pgfpathcurveto{\pgfpoint{56.946mm}{-86.970mm}}{\pgfpoint{56.938mm}{-87.225mm}}{\pgfpoint{56.906mm}{-87.337mm}}%
\pgfpathcurveto{\pgfpoint{56.715mm}{-87.982mm}}{\pgfpoint{56.124mm}{-88.357mm}}{\pgfpoint{55.439mm}{-88.278mm}}%
\pgfpathclose%
\pgfpathmoveto{\pgfpoint{70.255mm}{-88.238mm}}%
\pgfpathcurveto{\pgfpoint{69.545mm}{-88.046mm}}{\pgfpoint{69.131mm}{-87.329mm}}{\pgfpoint{69.314mm}{-86.619mm}}%
\pgfpathcurveto{\pgfpoint{69.338mm}{-86.523mm}}{\pgfpoint{69.553mm}{-86.069mm}}{\pgfpoint{69.825mm}{-85.550mm}}%
\pgfpathcurveto{\pgfpoint{70.606mm}{-84.035mm}}{\pgfpoint{70.989mm}{-83.301mm}}{\pgfpoint{71.411mm}{-82.480mm}}%
\pgfpathcurveto{\pgfpoint{71.635mm}{-82.049mm}}{\pgfpoint{71.914mm}{-81.515mm}}{\pgfpoint{72.033mm}{-81.284mm}}%
\pgfpathcurveto{\pgfpoint{72.153mm}{-81.061mm}}{\pgfpoint{72.440mm}{-80.510mm}}{\pgfpoint{72.671mm}{-80.064mm}}%
\pgfpathcurveto{\pgfpoint{72.903mm}{-79.617mm}}{\pgfpoint{73.254mm}{-78.939mm}}{\pgfpoint{73.461mm}{-78.549mm}}%
\pgfpathcurveto{\pgfpoint{73.660mm}{-78.158mm}}{\pgfpoint{73.963mm}{-77.560mm}}{\pgfpoint{74.139mm}{-77.209mm}}%
\pgfpathcurveto{\pgfpoint{74.514mm}{-76.467mm}}{\pgfpoint{74.593mm}{-76.364mm}}{\pgfpoint{74.968mm}{-76.180mm}}%
\pgfpathcurveto{\pgfpoint{75.526mm}{-75.909mm}}{\pgfpoint{76.212mm}{-76.013mm}}{\pgfpoint{76.635mm}{-76.443mm}}%
\pgfpathcurveto{\pgfpoint{76.786mm}{-76.603mm}}{\pgfpoint{76.786mm}{-76.595mm}}{\pgfpoint{77.671mm}{-78.309mm}}%
\pgfpathcurveto{\pgfpoint{77.879mm}{-78.708mm}}{\pgfpoint{78.182mm}{-79.282mm}}{\pgfpoint{78.333mm}{-79.569mm}}%
\pgfpathcurveto{\pgfpoint{78.493mm}{-79.864mm}}{\pgfpoint{78.772mm}{-80.399mm}}{\pgfpoint{78.963mm}{-80.766mm}}%
\pgfpathcurveto{\pgfpoint{79.147mm}{-81.132mm}}{\pgfpoint{79.482mm}{-81.778mm}}{\pgfpoint{79.705mm}{-82.201mm}}%
\pgfpathcurveto{\pgfpoint{80.526mm}{-83.772mm}}{\pgfpoint{81.866mm}{-86.411mm}}{\pgfpoint{81.906mm}{-86.547mm}}%
\pgfpathcurveto{\pgfpoint{82.257mm}{-87.727mm}}{\pgfpoint{81.069mm}{-88.676mm}}{\pgfpoint{79.944mm}{-88.110mm}}%
\pgfpathcurveto{\pgfpoint{79.617mm}{-87.943mm}}{\pgfpoint{79.490mm}{-87.775mm}}{\pgfpoint{79.027mm}{-86.882mm}}%
\pgfpathcurveto{\pgfpoint{78.485mm}{-85.821mm}}{\pgfpoint{78.509mm}{-85.861mm}}{\pgfpoint{78.381mm}{-85.805mm}}%
\pgfpathcurveto{\pgfpoint{78.254mm}{-85.750mm}}{\pgfpoint{73.110mm}{-85.726mm}}{\pgfpoint{72.911mm}{-85.781mm}}%
\pgfpathcurveto{\pgfpoint{72.743mm}{-85.821mm}}{\pgfpoint{72.775mm}{-85.781mm}}{\pgfpoint{72.081mm}{-87.145mm}}%
\pgfpathcurveto{\pgfpoint{71.587mm}{-88.126mm}}{\pgfpoint{71.053mm}{-88.445mm}}{\pgfpoint{70.255mm}{-88.238mm}}%
\pgfpathclose%
\pgfpathmoveto{\pgfpoint{77.097mm}{-83.278mm}}%
\pgfpathcurveto{\pgfpoint{77.265mm}{-83.214mm}}{\pgfpoint{77.257mm}{-83.174mm}}{\pgfpoint{76.730mm}{-82.081mm}}%
\pgfpathcurveto{\pgfpoint{76.507mm}{-81.627mm}}{\pgfpoint{76.292mm}{-81.172mm}}{\pgfpoint{76.252mm}{-81.077mm}}%
\pgfpathcurveto{\pgfpoint{76.212mm}{-80.981mm}}{\pgfpoint{76.085mm}{-80.718mm}}{\pgfpoint{75.981mm}{-80.494mm}}%
\pgfpathcurveto{\pgfpoint{75.622mm}{-79.745mm}}{\pgfpoint{75.646mm}{-79.729mm}}{\pgfpoint{75.024mm}{-81.021mm}}%
\pgfpathcurveto{\pgfpoint{74.003mm}{-83.134mm}}{\pgfpoint{73.987mm}{-83.174mm}}{\pgfpoint{74.115mm}{-83.262mm}}%
\pgfpathcurveto{\pgfpoint{74.187mm}{-83.317mm}}{\pgfpoint{76.970mm}{-83.325mm}}{\pgfpoint{77.097mm}{-83.278mm}}%
\pgfpathclose\pgfusepath{fill}%
}
\newcommand{\zuka@logo}[3]{\resizebox{!}{#2}{\color{#1}\begin{pgfpicture}\color{#1}#3\end{pgfpicture}}}
\newcommand{\zukamark}[2][zink]{\zuka@logo{#1}{#2}{\zuka@dibmark}}
\newcommand{\zukaver}[2][zink]{\zuka@logo{#1}{#2}{\zuka@dibmark\zuka@dibword}}
\newcommand{\zukahor}[2][zink]{\zuka@logo{#1}{#2}{\zuka@dibmark
  \begin{pgfscope}\pgftransformshift{\pgfpoint{54mm}{123.8mm}}\pgftransformscale{2}\zuka@dibword\end{pgfscope}}}


% ---------- color ----------
\definecolor{zink}{HTML}{0B1424}
\definecolor{zolive}{HTML}{173F73}   % azul tinta: títulos y rótulos
\definecolor{zmoss}{HTML}{4D7C0F}
\definecolor{zmute}{HTML}{56637A}
\definecolor{zline}{HTML}{CBD3DF}
\definecolor{zsoft}{HTML}{F1F5FA}
\definecolor{zpaper}{HTML}{F8FAFD}
\definecolor{zsun}{HTML}{B7791F}
\definecolor{zfire}{HTML}{D9541E}    % lo negativo o volteado, como en la app

% ---------- rótulo: barra de color, sin icono ----------
\renewcommand{\zuka@cab}[3]{{\color{#1}\rule[-.25ex]{1.1mm}{1.75ex}}\hspace{2mm}{\sffamily\footnotesize\bfseries\color{#1!85!black}#3}}
\tcbset{zbase/.style={enhanced,breakable,frame hidden,boxrule=0pt,arc=1.2mm,left=4.5mm,right=4.5mm,top=2.6mm,bottom=2.8mm,
  before skip=8pt,after skip=8pt,fontupper=\normalsize},
  zborde/.style={},
  zlisa/.style={enhanced,breakable,frame hidden,boxrule=0pt,sharp corners,colback=white,left=0mm,right=0mm,top=2.4mm,bottom=1mm,
    before skip=8pt,after skip=8pt,borderline north={.5pt}{0pt}{zline}},
  zmarco/.style={enhanced,breakable,colback=white,colframe=zolive,boxrule=.8pt,arc=1.2mm,left=4.5mm,right=4.5mm,top=2.8mm,bottom=3mm,
    before skip=9pt,after skip=9pt}}
\renewcommand{\zuka@cajaanalogia}[2]{%
  \begin{tcolorbox}[zbase,colback=zsoft]\zuka@cab{zolive}{idea}{La analogía}{\sffamily\footnotesize\color{zmute}\quad #1}\par\vspace{3pt}#2\end{tcolorbox}}
\renewenvironment{regla}[1][La regla]{\begin{tcolorbox}[zmarco]\zuka@cab{zolive}{regla}{#1}\par\vspace{4pt}}{\end{tcolorbox}}
\renewenvironment{ejemplo}[1][]{\stepcounter{ejemplo}%
  \begin{tcolorbox}[zlisa]\zuka@cab{zolive}{ejemplo}{Ejemplo \thesection.\arabic{ejemplo}}%
  \if\relax\detokenize{#1}\relax\else{\sffamily\footnotesize\color{zmute}\quad #1}\fi\par\vspace{3pt}}{\end{tcolorbox}}
\renewenvironment{trampa}{\begin{tcolorbox}[zbase,colback=zfire!6]\zuka@cab{zfire}{trampa}{Trampa}\par\vspace{2pt}}{\end{tcolorbox}}
\renewenvironment{atajo}{\begin{tcolorbox}[zbase,colback=zsoft]\zuka@cab{zolive}{atajo}{Atajo de examen}\par\vspace{2pt}}{\end{tcolorbox}}
\renewenvironment{profundiza}{\begin{tcolorbox}[zlisa]\zuka@cab{zmute}{lejos}{Para ir más lejos}\par\vspace{2pt}}{\end{tcolorbox}}
\newcommand{\zuka@titulo}[1]{\par\vspace{9pt}\noindent{\color{zlime}\rule[-.2ex]{7mm}{.9mm}}\hspace{2.5mm}%
  {\sffamily\fontsize{14}{17}\selectfont\fontseries{sb}\selectfont\color{zink}#1}\par\nopagebreak\vspace{3pt}}
\renewenvironment{reto}{\zuka@titulo{Ahora sí, un problema}%
  {\small\color{zmute}Ya tienes las ideas. Resuélvelo solo, escribe tu respuesta y después mira cómo se usó el concepto.\par}
  \begin{tcolorbox}[zmarco,breakable=false]}{\zuka@miresp\end{tcolorbox}}
\renewenvironment{retoresuelto}{\zuka@tapa{Tapa desde aquí hasta tener tu respuesta}%
  \begin{tcolorbox}[zbase,colback=zsoft]\zuka@cab{zmoss}{listo}{Así se usó el concepto}\par\vspace{3pt}}{\end{tcolorbox}}
\renewcommand{\zuka@cajamas}[1]{\setcounter{ejx}{0}\zuka@titulo{Ahora tú}%
  {\small\color{zmute}Resuélvelo aquí y compara con las claves del final.\par}#1\par\vspace{4pt}}
\renewcommand{\zuka@cajaexamen}[1]{\setcounter{forma}{0}\zuka@titulo{Así lo preguntan en el examen}#1\par\vspace{4pt}}
\renewcommand{\forma}[2]{\stepcounter{forma}\par\vspace{3pt}\noindent
  \begin{tabularx}{\linewidth}{@{}>{\leavevmode\sffamily\small\bfseries\color{zolive}}p{6mm}X@{}}
  \arabic{forma} & \leavevmode{\sffamily\small\bfseries\color{zink}#1}\par{\small #2}\end{tabularx}}
\renewcommand{\taller}{\zuka@ponemas\zuka@poneexamen\zuka@poneandamio\zuka@titulo{Taller \thesection}%
  {\sffamily\footnotesize\color{zmute}Resuélvelo todo antes de mirar las claves.\hfill
   \textbf{\color{zolive}Mis aciertos:} \rule[-1pt]{8mm}{.4pt} de \rule[-1pt]{8mm}{.4pt}}\par\vspace{2pt}}
\titleformat{\subsection}[hang]{\sffamily\fontsize{14}{17}\selectfont\fontseries{sb}\selectfont\color{zink}}{}{0pt}
  {{\color{zlime}\rule[-.2ex]{7mm}{.9mm}}\hspace{2.5mm}}
\titlespacing*{\subsection}{0pt}{11pt}{4pt}
% la bitácora solo queda en lectura: en razonamiento la reemplaza «Lo que construiste»
\let\zuka@bitacoravieja\bitacora
\renewcommand{\bitacora}{\ifnum\zuka@num<3\else\zuka@bitacoravieja\fi}

\raggedbottom
% el enunciado y su cuadrícula no se separan
\renewcommand{\ejx}[3]{\stepcounter{ejx}\par\vspace{4pt}\noindent\begin{minipage}{\linewidth}\setlength{\parskip}{4pt}%
  {\sffamily\small\bfseries\color{zolive}\Alph{ejx}\quad}{\sffamily\small\bfseries\color{zink}#1}\par
  #2\par\ifnum\zuka@num<3 \zuka@cuadricula{22mm}\else\zuka@renglones{2}\fi\end{minipage}\par
  \immediate\write\zuka@sol{\noexpand\solitem{Ahora tú \Alph{ejx}}{\unexpanded{#3}}}}
\titlecontents{chapter}[8mm]{\vspace{2pt}\sffamily\bfseries\color{zink}}
  {\contentslabel[\color{zolive}\two@digits{\thecontentslabel}]{8mm}}{\hspace*{-8mm}}{\hfill\color{zolive}\contentspage}
\titlecontents{section}[17mm]{\vspace{-1.6pt}\small\color{zink}}
  {\contentslabel[\color{zmute}\thecontentslabel]{9mm}}{}{\titlerule*[2.5mm]{\color{zline}.}\contentspage}

% ---------- cabecera de lección: número grande, título y ruta de la app ----------
\renewcommand{\zuka@cabezaleccion}[1]{%
  \noindent\begin{minipage}[b]{20mm}{\sffamily\fontsize{34}{34}\selectfont\firalight\color{zolive}\thesection}\end{minipage}%
  \begin{minipage}[b]{\dimexpr\linewidth-20mm}\raggedright\hyphenpenalty=10000{\sffamily\fontsize{22}{25}\selectfont\fontseries{sb}\selectfont\color{zink}#1}\end{minipage}%
  \par\vspace{4pt}\noindent{\color{zlime}\rule{20mm}{1mm}}{\color{zline}\rule[.35mm]{\dimexpr\linewidth-20mm}{.4pt}}}
\titleformat{\section}[block]{\sffamily\color{zink}}{}{0pt}{\zuka@cabezaleccion}
\titlespacing*{\section}{0pt}{0pt}{5pt}

\newcommand{\rutade}[2]{\expandafter\gdef\csname zrut@#1\endcsname{\zuka@ruta{#2}}}
\newcommand{\zuka@ruta}[1]{\zuka@np=0 \foreach \t in {#1}{\global\advance\zuka@np by 1 }%
  \par\vspace{4pt}\noindent\begin{tikzpicture}[x=1mm,y=1mm]
    \pgfmathsetmacro{\zw}{170/\the\zuka@np}
    \draw[zline,line width=1pt] (\zw/2,0) -- (170-\zw/2,0);
    \foreach \t [count=\k] in {#1}{%
      \ifnum\k=\zuka@np \filldraw[fill=zlime,draw=zlime,line width=1pt] (\k*\zw-\zw/2,0) circle (3);
      \else \filldraw[fill=white,draw=zolive,line width=1pt] (\k*\zw-\zw/2,0) circle (3);\fi
      \node[font=\sffamily\footnotesize\bfseries,text=zink] at (\k*\zw-\zw/2,0) {\k};
      \pgfmathsetmacro{\zl}{\zw-4}
      \node[anchor=north,text width=\zl mm,align=center,font=\sffamily\footnotesize,text=zmute] at (\k*\zw-\zw/2,-3.8) {\t};}
  \end{tikzpicture}\par\vspace{3pt}}
\newcommand{\zuka@poneruta}{\ifcsname zrut@\zuka@lid\endcsname\csname zrut@\zuka@lid\endcsname\fi}
\newcommand{\andamiode}[2]{\expandafter\gdef\csname zand@#1\endcsname{\zuka@andamio{#2}}}
\newcommand{\zuka@andamio}[1]{%
  \begin{tcolorbox}[zbase,colback=zlime!13]\zuka@cab{zmoss}{listo}{Lo que construiste}{\sffamily\footnotesize\color{zmute}\quad
  marca lo que podrías explicarle a otra persona; si algo queda sin marcar, vuelve a ese paso aquí o en la app}\par\vspace{3pt}
  \begin{itemize}[leftmargin=6mm,itemsep=2pt,label={\casilla{}}]\small #1\end{itemize}\end{tcolorbox}}
\newcommand{\zuka@poneandamio}{\ifcsname zand@\zuka@lid\endcsname\csname zand@\zuka@lid\endcsname
  \expandafter\global\expandafter\let\csname zand@\zuka@lid\endcsname\relax\fi}
\let\zuka@notaappvieja\notaapp
\renewcommand{\notaapp}[1]{\ifnum\zuka@num<3 \zuka@poneruta
    {\sffamily\footnotesize\color{zmute}\textbf{\color{zolive}En la app:} la misma ruta, para construirla con el dedo, y después los ejercicios.\par}%
  \else\zuka@notaappvieja{#1}\fi}
\InputIfFileExists{explora-rl}{}{}

% ---------- piezas de la lección por pasos ----------
\newcounter{zcap}\newcounter{zturno}
\renewcommand{\leccion}[3]{\global\zuka@evalfalse\setcounter{zcap}{0}\setcounter{zturno}{0}\zuka@leccionvieja{#1}{#2}{#3}\zuka@poneanalogia}
\newcommand{\capitulo}[1]{\stepcounter{zcap}\par\vspace{12pt}\noindent
  \tikz[baseline=(n.base)]{\node[circle,draw=zolive,line width=1pt,inner sep=0,minimum size=6.4mm,font=\sffamily\small\bfseries,text=zink](n){\arabic{zcap}};}%
  \hspace{2.5mm}{\sffamily\fontsize{15}{18}\selectfont\fontseries{sb}\selectfont\color{zink}#1}\par\nopagebreak\vspace{3pt}}
% \turno{enunciado (con su dibujo)}{respuesta para las claves}
\newcommand{\turno}[2]{\stepcounter{zturno}%
  \begin{tcolorbox}[zbase,breakable=false,colback=zlime!13]\zuka@cab{zmoss}{reto}{Tu turno}\quad #1\end{tcolorbox}%
  \immediate\write\zuka@sol{\noexpand\solitem{Tu turno \arabic{zturno}}{\unexpanded{#2}}}}
% \armar{\arma{título}{enunciado}{respuesta y el error que detecta}...}: evalúa la idea construyendo algo que cumpla una condición
\newcounter{zarma}
\newcommand{\armar}[1]{\setcounter{zarma}{0}\zuka@titulo{Arma la respuesta}%
  {\small\color{zmute}Como en los ejercicios de la app: no se pide un resultado, se pide construir algo que cumpla la condición.
   Escribe primero qué crees que pasa y después compruébalo con la cuenta o el dibujo.\par}#1\par}
\newcommand{\arma}[3]{\stepcounter{zarma}%
  \begin{tcolorbox}[zlisa,breakable=false]\zuka@cab{zolive}{taller}{\Alph{zarma}}{\sffamily\small\bfseries\color{zink}\quad #1}\par\vspace{2pt}#2\par\vspace{3pt}
  \zuka@cuadricula{20mm}\end{tcolorbox}%
  \immediate\write\zuka@sol{\noexpand\solitem{Arma la respuesta \Alph{zarma}}{\unexpanded{#3}}}}
\newcommand{\renglones}[1]{\zuka@renglones{#1}}
\newcommand{\cuenta}[1]{\par\vspace{1pt}\begin{center}\tikz\node[fill=zsoft,rounded corners=1.2mm,inner xsep=4mm,inner ysep=1.8mm]{$\displaystyle #1$};\end{center}\vspace{-2pt}}
\newcommand{\pie}[1]{{\par\centering\sffamily\footnotesize\color{zmute}#1\par}}

% ---------- dibujos: recta, barras, cien casillas, cuadrados de puntos ----------
\newcommand{\xficha}[2][zolive]{\filldraw[fill=#1,draw=white,line width=1pt] (#2,0) circle (2mm);}
\newcommand{\xsalto}[4][zolive]{\draw[#1,line width=1.1pt,-{Stealth[length=2.3mm]},shorten <=1.6mm,shorten >=1.6mm] (#2,.12) to[bend left=\ifdim#2pt<#3pt 45\else-45\fi] node[above,font=\sffamily\footnotesize\bfseries,text=#1]{#4} (#3,.12);}
\renewcommand{\recta}[3]{\begin{tikzpicture}[x=7.4mm,font=\sffamily\footnotesize]
  \draw[zline,line width=1.6pt,{Stealth[length=2mm]}-{Stealth[length=2mm]}] (#1-.9,0)--(#2+.9,0);
  \foreach \t in {#1,...,#2}{\draw[zmute,line width=.8pt] (\t,-.09)--(\t,.09); \node[below=1.4mm,text=zink] at (\t,0) {$\t$};}
  \draw[zink,line width=1.6pt] (0,-.14)--(0,.14);
  #3
\end{tikzpicture}}
\newcommand{\xrecta}[3]{\begin{center}\recta{#1}{#2}{#3}\end{center}}
% \xbarra[color]{y}{partes}{tomadas}{rótulo}: barra de 96 mm
\newcommand{\xbarra}[5][zlime]{\pgfmathsetmacro{\zp}{96/#3}
  \foreach \i in {1,...,#3}{\ifnum\i>#4 \filldraw[fill=white,draw=zmute,line width=.5pt] ({(\i-1)*\zp},#2) rectangle ++(\zp,8);
    \else \filldraw[fill=#1!75,draw=zmute,line width=.5pt] ({(\i-1)*\zp},#2) rectangle ++(\zp,8);\fi}
  \node[anchor=west,font=\small,text=zink] at (99,#2+4) {#5};}
\newenvironment{barras}{\begin{center}\begin{tikzpicture}[x=1mm,y=1mm]}{\end{tikzpicture}\end{center}}
% \xcien{x}{pintadas}{segunda cantidad pintada en naranja}{pie}
\newcommand{\xcien}[4]{\foreach \j in {0,...,9}{\foreach \i in {0,...,9}{\pgfmathtruncatemacro{\n}{\j*10+\i}
    \ifnum\n<#2 \def\zc{zlime!80}\else\pgfmathtruncatemacro{\m}{#2+#3}\ifnum\n<\m \def\zc{zfire!55}\else\def\zc{white}\fi\fi
    \filldraw[fill=\zc,draw=zline,line width=.4pt] (#1+\i*3.2,-\j*3.2) rectangle ++(3.2,-3.2);}}
  \draw[zolive,line width=.8pt] (#1,0) rectangle ++(32,-32);
  \node[anchor=north,align=center,font=\sffamily\footnotesize,text=zmute] at (#1+16,-33.5) {#4};}
\newenvironment{casillas}{\begin{center}\begin{tikzpicture}[x=1mm,y=1mm]}{\end{tikzpicture}\end{center}}
% \xcuad{x}{lado}{tamaño de casilla}{pie}
\newcommand{\xcuad}[4]{\foreach \i in {1,...,#2}{\foreach \j in {1,...,#2}{%
    \filldraw[fill=zlime!70,draw=white,line width=.8pt] ({#1+(\i-1)*#3},{(\j-1)*#3}) rectangle ++(#3,#3);}}
  \node[anchor=north,align=center,font=\sffamily\footnotesize,text=zmute] at ({#1+#2*#3/2},-1.5) {#4};}

% ---------- encabezado, pie y pestaña ----------
\renewcommand{\zuka@pestana}[1]{\ifnum\value{mundo}>0
  \begin{tikzpicture}[remember picture,overlay]
    \pgfmathsetmacro{\zty}{-26-mod(\value{mundo}-1,12)*15}
    \ifnum#1=1
      \fill[zlime] ([xshift=-5mm,yshift=\zty mm]current page.north east) rectangle ++(6mm,-11mm);
      \node[font=\sffamily\scriptsize\bfseries,text=zink] at ([xshift=-2.5mm,yshift=\zty mm-5.5mm]current page.north east) {\themundo};
    \else
      \fill[zlime] ([xshift=5mm,yshift=\zty mm]current page.north west) rectangle ++(-6mm,-11mm);
      \node[font=\sffamily\scriptsize\bfseries,text=zink] at ([xshift=2.5mm,yshift=\zty mm-5.5mm]current page.north west) {\themundo};
    \fi
  \end{tikzpicture}\fi}
\fancyhf{}
\fancyhead[L]{\zuka@pestana{\ifodd\value{page}1\else0\fi}\sffamily\footnotesize\color{zmute}\zuka@mundoactual}
\fancyhead[R]{\sffamily\footnotesize\color{zmute}\zuka@comp\enspace·\enspace Libro \zuka@num}
\fancyfoot[L]{\raisebox{-.6mm}{\zukahor[zolive]{4.6mm}}\sffamily\footnotesize\color{zmute}\enspace Preuniversitario UdeA}
\fancyfoot[R]{\sffamily\small\bfseries\color{zolive}\thepage}
\fancypagestyle{plain}{\fancyhf{}\fancyfoot[R]{\sffamily\small\bfseries\color{zolive}\thepage}}

% ---------- apertura de mundo: clara, con el arte a la derecha ----------
\renewcommand{\mundo}[4]{%
  \cleardoublepage
  \stepcounter{mundo}\setcounter{section}{0}\global\zuka@evalfalse
  \renewcommand{\zuka@mundoactual}{Mundo \themundo\ · #1}%
  \thispagestyle{plain}%
  \phantomsection\addcontentsline{toc}{chapter}{\protect\numberline{\themundo}#1}%
  \begin{tikzpicture}[remember picture,overlay]
    \fill[zsoft] (current page.north west) rectangle ([yshift=-92mm]current page.north east);
    \begin{scope}
      \clip ([xshift=-72mm]current page.north east) rectangle ([yshift=-92mm]current page.north east);
      \pgfmathsetmacro{\zdy}{46+mod(\themundo*41,90)}
      \node[anchor=north east,inner sep=0] at ([xshift=26mm,yshift=\zdy mm]current page.north east)
        {\includegraphics[width=150mm]{arte/portada\zuka@num.jpg}};
    \end{scope}
    \fill[zlime] ([yshift=-92mm]current page.north west) rectangle ++(20mm,-1.6mm);
    \fill[zlime] ([xshift=20mm,yshift=-92mm]current page.north west) rectangle ++(22mm,-1.6mm);
    \node[anchor=north west,inner sep=0] at ([xshift=20mm,yshift=-16mm]current page.north west) {\zukahor[zink]{9mm}};
    \node[anchor=north west,inner sep=0,font=\sffamily\small,text=zmute]
      at ([xshift=20mm,yshift=-38mm]current page.north west) {\zuka@comp\enspace·\enspace\textbf{\color{zolive}Mundo \themundo}};
    \node[anchor=north west,inner sep=0,text width=112mm,align=flush left,font=\sffamily\fontsize{31}{35}\selectfont\fontseries{sb}\selectfont,text=zink]
      at ([xshift=20mm,yshift=-46mm]current page.north west) {#1};
  \end{tikzpicture}%
  \par\vspace*{80mm}%
  {\large #3\par}\vspace{7mm}%
  \noindent{\sffamily\fontsize{14}{17}\selectfont\fontseries{sb}\selectfont\color{zink}Las lecciones}\par\vspace{2pt}
  {\color{zline}\rule{\linewidth}{.5pt}}\par
  {\begin{itemize}[leftmargin=0mm,itemsep=5pt,label={}]#4\end{itemize}}
  \vspace{7mm}
  \noindent{\sffamily\fontsize{14}{17}\selectfont\fontseries{sb}\selectfont\color{zink}Cómo vas a trabajar}\par\vspace{2pt}
  {\color{zline}\rule{\linewidth}{.5pt}}\par\vspace{2pt}
  {\small\ifnum\zuka@num<3
    Cada lección es una \textbf{ruta}: los mismos pasos que construyes en la app. En cada paso \textbf{primero entiendes la
    idea} y \textbf{después la usas}; al final viene el problema y ves cómo se usó el concepto. Cierra con «Lo que
    construiste», el taller y las claves.
  \else
    En cada lección \textbf{primero entiendes la técnica} (analogía, idea y ejemplos comentados) y \textbf{después la usas}:
    resuelves el reto, ves cómo se aplicó, practicas en la app y cierras con el taller y las claves.
  \fi\par}
  \immediate\write\zuka@sol{\noexpand\solmun{Mundo \themundo\ · \unexpanded{#1}}}%
}

% ---------- evaluación, portada y contraportada ----------
\renewcommand{\zuka@cabeceravieja}[3]{%
  \noindent{\sffamily\footnotesize\bfseries\color{zolive}#1}\hfill{\sffamily\footnotesize\color{zmute}\begin{tabular}[t]{@{}r@{}}#3\end{tabular}}\par\vspace{1pt}
  {\sffamily\fontsize{22}{25}\selectfont\fontseries{sb}\selectfont\color{zink}#2\par}\vspace{4pt}
  \noindent{\color{zlime}\rule{20mm}{1mm}}{\color{zline}\rule[.35mm]{\dimexpr\linewidth-20mm}{.4pt}}\par}
\renewcommand{\portada}{%
  \thispagestyle{vacia}%
  \begin{tikzpicture}[remember picture,overlay]
    \fill[white] (current page.south west) rectangle (current page.north east);
    \begin{scope}
      \clip ([yshift=34mm]current page.south west) rectangle ([yshift=-138mm]current page.north east);
      \node[inner sep=0] at ([yshift=-22mm]current page.center) {\includegraphics[width=\paperwidth]{arte/portada\zuka@num.jpg}};
    \end{scope}
    \node[anchor=north west,inner sep=0] at ([xshift=20mm,yshift=-18mm]current page.north west) {\zukahor[zink]{15mm}};
    \node[anchor=north east,inner sep=0,align=right,font=\sffamily\small,text=zmute] at ([xshift=-20mm,yshift=-21mm]current page.north east)
      {Preuniversitario UdeA\\\textbf{\color{zolive}Libro \zuka@num\ de 4}};
    \fill[zlime] ([xshift=20mm,yshift=-50mm]current page.north west) rectangle ++(26mm,-2mm);
    \node[anchor=north west,inner sep=0,text width=170mm,align=flush left]
      at ([xshift=20mm,yshift=-60mm]current page.north west)
      {{\sffamily\fontsize{20}{24}\selectfont\firalight\color{zolive}\zuka@comp\par}\vspace{4mm}%
       {\sffamily\fontsize{40}{44}\selectfont\fontseries{sb}\selectfont\color{zink}\zuka@tit\par}\vspace{6mm}%
       \parbox{122mm}{\raggedright\sffamily\small\color{zmute}\zuka@sub}};
    \node[anchor=north west,inner sep=0,text width=170mm,align=flush left,font=\sffamily\footnotesize,text=zmute]
      at ([xshift=20mm,yshift=27mm]current page.south west) {\zuka@mundos};
  \end{tikzpicture}%
  \null\clearpage}
\renewcommand{\contraportada}{%
  \cleardoublepage\thispagestyle{vacia}%
  \begin{tikzpicture}[remember picture,overlay]
    \fill[white] (current page.south west) rectangle (current page.north east);
    \begin{scope}
      \clip (current page.north west) rectangle ([yshift=-96mm]current page.north east);
      \node[anchor=north,inner sep=0] at ([yshift=62mm]current page.north) {\includegraphics[width=\paperwidth]{arte/portada\zuka@num.jpg}};
    \end{scope}
    \fill[zlime] ([xshift=20mm,yshift=-114mm]current page.north west) rectangle ++(26mm,-2mm);
    \node[anchor=north west,inner sep=0,text width=150mm,align=flush left,font=\sffamily\fontsize{24}{29}\selectfont\fontseries{sb}\selectfont,text=zink]
      at ([xshift=20mm,yshift=-122mm]current page.north west) {Primero la idea.\\Después, úsala.};
    \node[anchor=north west,inner sep=0,text width=140mm,align=flush left,text=zink]
      at ([xshift=20mm,yshift=-152mm]current page.north west)
      {Cuatro libros y una aplicación para preparar la prueba de admisión de la Universidad de Antioquia: dos de
      Razonamiento lógico y dos de Comprensión lectora. Cada lección del libro sigue el mismo camino que su lección en Zuka
      y va más hondo: analogías, ejemplos resueltos paso a paso, trampas, atajos, talleres y, al cerrar cada mundo, una
      evaluación como la del examen.};
    \foreach \n/\t in {1/Número y forma,2/{Orden, lógica y azar},3/Leer lo que dice,4/Leer lo que quiere decir}{%
      \node[anchor=west,inner sep=0,font=\sffamily\small,text=zink] at ([xshift=20mm,yshift=54mm-\n*7mm]current page.south west)
        {{\color{zolive}\bfseries Libro \n}\quad \t};}
    \node[anchor=south east,inner sep=0] at ([xshift=-20mm,yshift=24mm]current page.south east) {\zukaver[zink]{34mm}};
  \end{tikzpicture}}
